\subsection{Replacing a minimal ideal by the states of an act}

\begin{theorem}[Constant minimal-ideal act compression]
\label{C:constant-minimal-act-compression}
Let $M$ be a finite monoid with proper minimal ideal $K$, and let
$I$ be the set of its minimal right ideals (the row set of its
Rees matrix presentation).  Let $S$ be a nonempty finite right
$M$-set.  Suppose that every $k\in K$ acts on $S$ as a constant
map.  There is an explicitly defined finite monoid $N$ with
\[
 |N|=|M|-|K|+|I|\,|S|
 \quad\hbox{and}\quad
 BN\simeq B\mathcal C(M,S).
\]
Here $\mathcal C(M,S)$ has objects $a,b$, endomorphism monoids
$M,\{1_b\}$, hom-set $\Hom(a,b)=S$, and no arrows from $b$ to $a$.
The conclusion allows nontrivial maximal groups in $K$.
\end{theorem}
\begin{proof}
Write $c_k\in S$ for the value of the constant action of $k$.
Right actions give $c_{kl}=c_l$ for $k,l\in K$.  In a completely
simple Rees matrix ideal, left multiplication by elements of $K$
varies the row and group coordinate arbitrarily within any fixed
column.  Consequently $c_k$ depends only on the column of $k$.
Let $\rho:K\to I$ be its row map.  The map
\[
 \eta:K\longrightarrow I\times S,
 \qquad \eta(k)=(\rho(k),c_k)
\]
is a homomorphism to the rectangular band
$(i,s)(j,t)=(i,t)$.  Left multiplication by $m\in M$ induces a
map $\sigma_m:I\to I$.  The constant-action identities imply
\[
 \eta(mk)=(\sigma_m\rho(k),c_k),\qquad
 \eta(km)=(\rho(k),c_km).                         \tag{\ref{C:constant-minimal-act-compression}.1}
\]

Set $N=(M\setminus K)\sqcup(I\times S)$.  Products of two old
elements are their product in $M$ if it lies outside $K$, and its
$\eta$-image otherwise.  Define all other products by
\[
 m(i,s)=(\sigma_m(i),s),\qquad
 (i,s)m=(i,sm),\qquad
 (i,s)(j,t)=(i,t).
\]
Equations (\ref{C:constant-minimal-act-compression}.1) verify every
associativity equation that meets the new band; the remaining
equations are those in $M$.  The old identity remains an identity.
Thus $N$ is a monoid and the identity outside $K$, together with
$\eta$ on $K$, defines a unital homomorphism $q:M\to N$.

Fix $i_0\in I,s_0\in S$ and put $e=(i_0,s_0)\in N$.
Let $D$ be the full subcategory of the Karoubi envelope of $N$
on the idempotents $1,e$.  Its four hom-sets are
\begin{align*}
 \End_D(1)&=N, & \End_D(e)&=\{e\},\\
 \Hom_D(1,e)&=eN=\{(i_0,s):s\in S\}, &
 \Hom_D(e,1)&=Ne=\{(i,s_0):i\in I\}.
\end{align*}
The object $e$ is a retract of $1$, so $BD\simeq BN$.
There is a functor $F:\mathcal C(M,S)\to D$ sending $a$ to $1$,
$b$ to $e$, $m$ to $q(m)$, and $s$ to $(i_0,s)$.
We verify Quillen's Theorem A directly for this functor.

The comma category $(F\downarrow e)$ has terminal object
$(b,e)$.  Indeed an object over $a$ is an element $(i_0,s)$,
and the unique arrow to $(b,e)$ is its corresponding $s\in S$.
For the other comma category $E=(F\downarrow1)$, the objects
over $a$ are the elements $n\in N$, and the objects over $b$
are symbols $*_i$, one for each $(i,s_0)\in Ne$.  Its arrows
among the old objects are
\[
 n\xrightarrow{\ m\ }n'\quad\Longleftrightarrow\quad n'q(m)=n.
                                                               \tag{\ref{C:constant-minimal-act-compression}.2}
\]
There is one arrow $n\to *_i$ exactly when $n=(i,s)$ for some
$s\in S$, and there are no other nonidentity arrows incident
from a new object.

For a right $M$-set $X$, let $A(X)$ denote the opposite of its
action category: its arrows are $x\to y$ labelled $m$ with
$ym=x$.  Choosing the final object and the arrow labels gives
a natural bijection $N_qA(X)=X\times M^q$.  In particular this
nerve construction preserves colimits of right $M$-sets and
preserves monomorphisms.  As right $M$-sets the construction of
$N$ is the pushout
\[
 N=M\sqcup_K(I\times S).                         \tag{\ref{C:constant-minimal-act-compression}.3}
\]
The full row category $A(\{i\}\times S)$ is a sieve in $A(N)$:
an arrow with target in that row has its source in the same row.
The description (\ref{C:constant-minimal-act-compression}.2)
therefore identifies the nerve of $E$ with the nerve of $A(N)$
with a terminal cone attached to each row nerve.  There are no
additional mixed simplices: a chain ending at $*_i$ has all its
previous objects in row $i$.  Using
(\ref{C:constant-minimal-act-compression}.3) gives the actual
pushout of realizations
\[
 BE=BA(M)\ \sqcup_{BA(K)}\
       \coprod_{i\in I}\operatorname{Cone}(BA(S)).                \tag{\ref{C:constant-minimal-act-compression}.4}
\]
The map $BA(K)\hookrightarrow BA(M)$ is a simplicial
subcomplex inclusion and hence a cofibration.  Thus this ordinary
pushout computes its homotopy pushout.  Contracting each cone to
its vertex consequently gives
\[
 BE\simeq BA(M)\sqcup^h_{BA(K)} I.
\]
The category $A(M)$ has terminal object $1$.  Also $K$ is the
disjoint union of the minimal right ideals $e_iM$, with one
idempotent $e_i$ for each $i\in I$.  Each $BA(e_iM)$ is
contractible.  More explicitly, the arrows
$x\to e_i$ labelled $x$ form a natural transformation from
the identity of $A(e_iM)$ to the constant functor at $e_i$:
for an arrow $x\xrightarrow{m}y$ its naturality equation is
$ym=x$.  The row projection $BA(K)\to I$ is therefore a
homotopy equivalence.  It follows that $BE$ is contractible.

Both comma categories of $F$ are contractible, so Quillen's
Theorem A gives $B\mathcal C(M,S)\simeq BD\simeq BN$.
\end{proof}

\begin{corollary}[A numerical obstruction from the monoid minimum]
\label{C:constant-minimal-act-size}
Under the hypotheses of
Theorem~\ref{C:constant-minimal-act-compression}, if
$B\mathcal C(M,S)\simeq\mathbb CP^2$, then
\[
 |M|-|K|+|I|\,|S|\geq23.
\]
\end{corollary}
\begin{proof}
Apply the exact monoid minimum in Theorem~\ref{B:CP2-minimum23}
to the monoid just constructed.
\end{proof}

\begin{remark}[The constant-action hypothesis is essential to the formula]
\label{C:constant-minimal-act-scope}
The theorem does not infer constant action from the number of
states.  In particular a nontrivial maximal group in $K$ can act
nontrivially on a small orbit.  Such an action must be treated
separately before applying this compression.  If $K=M$ is a
group and its identity acts constantly, $S$ is a singleton and
the act cone is contractible; this degenerate case needs no
replacement construction.
\end{remark}

\begin{lemma}[The constant-action quotient detects degree two]
\label{C:constant-action-quotient-H2}
Let $M$ be any finite monoid with minimal ideal $K$, and let $S$
be a nonempty finite right $M$-set.  Let $\bar S$ be the quotient
of $S$ by the equivalence relation generated by
\[
 sk\sim s'k\qquad(s,s'\in S, k\in K).
\]
Then $\bar S$ is a right $M$-set on which every element of $K$
acts constantly.  The induced functor of act cones gives a
surjection
\[
 H^2(B\mathcal C(M,\bar S);\mathbb Q)
 \longrightarrow H^2(B\mathcal C(M,S);\mathbb Q).
\]
If $K$ is a rectangular band, the same surjection holds with
integral coefficients.
In particular, if $|S|\leq3$ and some element of $K$ does not
act constantly on $S$, all degree-two rational products on
$B\mathcal C(M,S)$ vanish.  No restriction on the maximal group
of $K$ is imposed.
\end{lemma}
\begin{proof}
The relation is $M$-invariant because $km\in K$ for $k\in K$,
$m\in M$.  Put $R=\mathbb QM$, $L=\widetilde{\mathbb QS}$,
and $\bar L=\widetilde{\mathbb Q\bar S}$.  Its relation span
is exactly
\[
 \ker(L\longrightarrow\bar L)=LK:
 \qquad ([s]-[s'])k=[sk]-[s'k].
\]
Here taking the transitive closure introduces no additional linear
relations, since the difference between two points in an equivalence
class is the sum of differences along a generating chain.

Choose a minimal idempotent and average its maximal group, as in
Theorem~\ref{C:two-row-act-cups}.  This gives a projective left
$R$-module $P\cong\mathbb QI$ and an exact sequence
\[
 0\longrightarrow V\longrightarrow P\longrightarrow\mathbb Q
       \longrightarrow0,
 \qquad V=\widetilde{\mathbb QI},\quad KV=0.
\]
For the rational dual $D$, every map $V\to DL$ takes values in
the annihilator of $LK$: for $v\in V$, $\ell\in L$ and $k\in K$,
equivariance gives $f(v)(\ell k)=(kf(v))(\ell)=f(kv)(\ell)=0$.
Thus the inclusion $D\bar L\hookrightarrow DL$ induces an
isomorphism
\[
 \Hom_R(V,D\bar L)\xrightarrow{\ \cong\ }\Hom_R(V,DL).
\]
The natural act-cone bar and the above projective presentation
identify its degree-two cohomology with
\[
 \operatorname{Ext}^1_R(\mathbb Q,DL)
  =\operatorname{coker}\bigl(\Hom_R(P,DL)
                         \longrightarrow\Hom_R(V,DL)\bigr).
\]
Applying the same formula to $\bar L$ proves the claimed
surjection: the numerator is unchanged and the denominator can
only increase on passing to $DL$.

For a minimal band the same argument works over $\mathbb Z$.
Here $P=\mathbb ZI$ is projective without averaging, and
$L/LK=\widetilde{\mathbb Z\bar S}$ is a free abelian group:
the subgroup $LK$ is exactly the lattice generated by differences
inside equivalence classes.  Its integral dual is therefore the
annihilator of $LK$, and the same Hom-isomorphism and connecting
cokernels prove the integral assertion.

If a $K$-actor is nonconstant and $|S|\leq3$, the quotient has
at most two points.  Every action on a two-point set is affine
on $\{0,1\}$ with slope in $\{0,1,-1\}$, and a one-point act
has contractible cone.  Lemma~\ref{C:finite-line-act-cups} makes
all positive-degree products of the quotient cone zero.
The cohomology surjection and cup naturality prove the last claim.
\end{proof}

\begin{lemma}[A three-by-six band with five outside elements]
\label{C:three-by-six-five-outside-CP2}
A finite monoid with rectangular-band minimal ideal of shape
$3\times6$ or $6\times3$, and exactly five elements outside that
ideal, cannot have classifying space $\mathbb CP^2$.
\end{lemma}
\begin{proof}
After opposition the band--Tor modules have dimensions
$\dim V=2$, $\dim W=5$, over the five-dimensional contracted
monoid algebra $A$, with Tor dimensions $1,0,1,0,\ldots$.
A local $A$ is excluded by Proposition~\ref{B:local-quotient-CP2}.
The nonlocal classification in
Proposition~\ref{B:five-nonlocal-exceptions} leaves just the line
and Nakayama algebras of Lemma~\ref{B:five-exceptional-targets}.

For the zero-composite line algebra, its displayed resolutions
and $\dim V=2$, $\operatorname{Tor}_1=0$, $\dim\operatorname{Tor}_2=1$
force
\[
 V=S_1\oplus S_3,\qquad
 W=T_1\oplus T_3^{\oplus s}\oplus(e_1A)^{\oplus p}
                           \oplus(e_2A)^{\oplus q}.
\]
The dimension and tensor constraints are
$s+2p+2q=4$ and $1+s+p=1$.  Hence $s=p=0,q=2$.
Only the $T_1$ summand contributes to the balanced tensor product,
so the unique degree-two form has matrix rank one.  Its square
vanishes by Lemma~\ref{B:rank-one-square}.

For the Nakayama algebra, the full Tor table in the same lemma
forces $V=L$ and
$W=T_2\oplus(e_1A)^{\oplus p}\oplus(e_2A)^{\oplus q}$.
The dimension equation is $3p+2q=4$, with unique solution
$p=0,q=2$.  Its tensor dimension is $p+q=2$, contradicting the
required value one.  These cases are exhaustive and prove the claim.
\end{proof}

\begin{theorem}[The exact three-state act-cone minimum]
\label{C:three-state-cone-minimum22}
If a three-state act cone $\mathcal C(M,S)$ has classifying space
$\mathbb CP^2$, then it has at least twenty-two arrows.  This bound
is attained by the two models of Theorem~\ref{C:CP2-upper22}.
\end{theorem}
\begin{proof}
Suppose $m=|M|\leq17$, so the cone has at most twenty-one arrows.
Its nonzero degree-two square requires $H^2(BM;\mathbb Q)\ne0$;
otherwise every degree-two class is relative to $BM$ and products
vanish.  Write $|K|=rsh$, where $r,s$ are the numbers of rows and
columns and $h$ is the maximal-group order.  The rational balanced
presentation requires $r,s\geq2$, and
Theorem~\ref{C:two-row-act-cups} requires $r\geq3$.

Lemma~\ref{C:constant-action-quotient-H2} now forces every element
of $K$ to act constantly.  The compression is therefore an actual
homotopy-equivalent monoid, of order
\[
 |N|=m-rsh+3r.
\]
If $h\geq2$, then $sh\geq4$, and $|N|\leq m-r<23$,
contradicting Theorem~\ref{B:CP2-minimum23}.  If $h=1$ and
$s\geq3$, then $|N|\leq m<23$, with the same contradiction.
It remains that $h=1,s=2$.  Put $q=m-2r$, the number of elements
outside the minimal band.  Compression leaves precisely $q$
elements outside its new $r\times3$ band.  The exact monoid
minimum and Corollary~\ref{B:five-outside-minimal-band} require
\[
 m+r\geq23,\qquad q=m-2r\geq5.
\]
Together with $m\leq17$ these inequalities force
$m=17$, $r=6$, $q=5$.  The compressed monoid then has a
$6\times3$ band and five outside elements, excluded by
Lemma~\ref{C:three-by-six-five-outside-CP2}.
This proves the lower bound.  The models of
Theorem~\ref{C:CP2-upper22} have an eighteen-element endomorphism
monoid, three states and one sink identity, and give equality.
\end{proof}

\begin{lemma}[At most four states over a minimal band]
\label{C:four-state-band-dichotomy}
Let $M$ have rectangular-band minimal ideal $K=I\times J$, and
let $S$ be a nonempty finite right $M$-set with $|S|\leq4$.
If $H_1(B\mathcal C(M,S);\mathbb Q)=0$, then either every
degree-two rational cup product on this act cone is zero, or
every element of $K$ acts constantly on $S$.
\end{lemma}
\begin{proof}
First suppose that some $z\in S$ is fixed by every element of
$K$.  Adjoining a fresh identity to $K$ gives a unital monoid
$K^+\to M$ and a functor $K^+\to S\rtimes M$, taking its object
to $z$.  By the natural balanced-tensor presentation in
Lemma~\ref{lem:B-band-Tor}, the induced map
\[
 H_2(BK^+;\mathbb Q)
   =\widetilde{\mathbb QJ}\otimes_{\mathbb Q}
      \widetilde{\mathbb QI}
   \longrightarrow H_2(BM;\mathbb Q)
\]
is surjective: passing from $K^+$ to $M$ only adds balancing
relations.  Hence $H^2(BM;\mathbb Q)\to H^2(B(S\rtimes M);\mathbb Q)$
is injective.  The cofiber sequence for the act cone implies that
every class in its $H^2$ lifts to relative cohomology with respect
to $BM$.  The relative quotient is a suspension of the edge
space, so its positive-degree cup products vanish.  Their images
give the asserted zero products on the act cone.

Assume now that no point is fixed by all of $K$.  The set $SK$
is the disjoint union of its right $K$-orbits.  Each such orbit
is $M$-invariant, since $(xk)m=x(km)$ and $km\in K$.
It has at least two points, since a singleton orbit would be a
$K$-fixed point.  If there are at least two orbits, $|S|\leq4$
forces exactly two two-point orbits and $S=SK$.  The action
category is then disconnected.  This is impossible: the
degree-zero and degree-one part of the act-cone exact sequence,
using $H_1(BM;\mathbb Q)=0$, identifies
$H_1(B\mathcal C(M,S);\mathbb Q)$ with the reduced $H_0$ of
that action category.

There is therefore a single right $K$-orbit in $SK$.
Every $k\in K$ acts constantly on this orbit.  Indeed choose a
point $y=xk_0$ in the orbit and write every point as
$y\ell=xk_0\ell$.  The band product $k_0\ell k$ has the row
of $k_0$ and the column of $k$, independently of $\ell$.
Thus $(y\ell)k$ is independent of the point $y\ell$.
Finally, for an arbitrary $x\in S$, idempotence gives
$xk=(xk)k$.  The point $xk$ lies in $SK$, so this is the same
constant value for every $x$.  Thus every $k$ acts constantly
on all of $S$.
\end{proof}

\begin{corollary}[Four-state cones through seventeen arrows]
\label{C:four-state-cones-seventeen}
If $|S|=4$ and $|M|\leq12$, then
$B\mathcal C(M,S)\not\simeq\mathbb CP^2$.
Thus no four-state act cone with at most seventeen arrows
models $\mathbb CP^2$.  Also, if $|S|=3$, $|M|\leq15$, and
the minimal ideal of $M$ is a band, the same exclusion holds.
\end{corollary}
\begin{proof}
A possible model must have $H^2(BM;\mathbb Q)\ne0$, since
otherwise every degree-two cone class is relative and has zero
products.  Write the minimal-ideal shape as $I\times H\times J$,
with $r=|I|,s=|J|$.  The rational group-collapse and balanced-tensor
formulas of Lemmas~\ref{lem:B-collapse-minimal-group} and
\ref{lem:B-band-Tor} require $r,s\geq2$.
Theorem~\ref{C:two-row-act-cups} further requires $r\geq3$.
When $|M|\leq12$, a nontrivial $H$ would already require
$rs|H|+1\geq13$ elements.  Thus $K$ is a band.

Lemma~\ref{C:four-state-band-dichotomy} now forces constant
$K$-action.  For four states the compressed monoid has order
\[
 |N|=m-rs+4r\leq m+2r
 \leq m+2\left\lfloor\frac{m-1}{2}\right\rfloor\leq22
 \qquad(m=|M|\leq12).
\]
This contradicts Corollary~\ref{C:constant-minimal-act-size}.
For three states and a band ideal, the same argument gives
\[
 |N|=m-rs+3r\leq m+r
 \leq m+\left\lfloor\frac{m-1}{2}\right\rfloor\leq22
 \qquad(m\leq15).
\]
The band qualification in this last assertion is explicit;
no inference of constant action is made for a nontrivial
minimal group.
\end{proof}

\begin{theorem}[A three-object triangle as a band extension]
\label{C:triangle-band-compression}
Let $M$ be a finite monoid, $S$ a nonempty finite left $M$-set,
and $T$ a nonempty finite right $M$-set.  Let $C$ have objects
$0,1,2$, endomorphism monoids $\{1_0\},M,\{1_2\}$, and
\[
 \Hom_C(0,1)=S,\qquad \Hom_C(1,2)=T,\qquad
 \Hom_C(0,2)=\{w\},
\]
with no other arrows between distinct objects.  Every composite
$t s$ equals $w$.  Then $BC\simeq BN$, where
\[
 N=M\sqcup(S\times T),\qquad |N|=|M|+|S|\,|T|,
\]
and multiplication extends that of $M$ by
\[
 (s,t)(s',t')=(s,t'),\qquad
 m(s,t)=(ms,t),\qquad (s,t)m=(s,tm).
\]
\end{theorem}
\begin{proof}
The two action laws verify associativity of $N$ directly.
Choose $s_0\in S,t_0\in T$ and put $e=(s_0,t_0)$.
In the full subcategory $D\subset\operatorname{Kar}(N)$ on
$1,e$, the corner $eNe$ is a singleton,
$\Hom_D(e,1)=Ne\cong S$, and $\Hom_D(1,e)=eN\cong T$.
Define $F:C\to D$ by sending $0,2$ to $e$ and $1$ to $1$,
using the inclusion $M\hookrightarrow N$, sending
$s\mapsto(s,t_0)$ and $t\mapsto(s_0,t)$, and sending $w$ to $e$.
The composite of the latter two band elements is $e$, so this
is a functor.  Again $BD\simeq BN$ because $e$ is a retract of $1$.

The comma category $(F\downarrow e)$ has terminal object
$(2,e)$.  To analyze $E=(F\downarrow1)$, denote its objects
over $0$ and $2$ by $0_s,2_s$ for $s\in S$, and its objects
over $1$ by $1_n$ for $n\in N$.  The arrows among the latter
objects are labelled $m\in M$ with $n'm=n$, as in the previous
proof.  An arrow $0_s\to1_n$ is labelled $a\in S$ with
$n a=s$, where $M$ acts on $S$ by the given action and
$(u,t)\in S\times T$ acts constantly with value $u$.
An arrow $1_n\to2_s$ exists precisely when $n=(s,t)$;
it is then uniquely labelled $t$.  Finally there is one arrow
$0_s\to2_{s'}$ if and only if $s=s'$, labelled $w$.

Let $E_0$ be the full subcategory on all $0_s$ and the
$1_m$ with $m\in M$.  It has terminal object $1_1$:
the unique arrow $1_m\to1_1$ is labelled $m$, and the unique
arrow $0_s\to1_1$ is labelled $s$.
Let $E_1$ be the full subcategory on all $0_s,2_s$ and all
$1_{(s,t)}$.  It is the disjoint union, over $s\in S$, of
categories with terminal object $2_s$.
Their intersection is the discrete category of the objects $0_s$.

There are no arrows between an old label $1_m$ and a band label
$1_{(s,t)}$: right multiplication by elements of $M$ preserves
the two summands of $N$.  The shared objects $0_s$ have no
nonidentity incoming arrows.  Every composable chain in $E$
therefore lies entirely in $E_0$ or entirely in $E_1$, and its
nerve is the actual pushout
\[
 NE=NE_0\sqcup_S NE_1.
\]
Both inclusions of $S$ are simplicial cofibrations.  Moreover
$S\to BE_1$, selecting $0_s$ in its contractible component, is
a homotopy equivalence.  The pushout is consequently homotopy
equivalent to $BE_0$, which is contractible.  Quillen's Theorem A
now gives $BC\simeq BD\simeq BN$.
\end{proof}

\begin{corollary}\label{C:triangle-band-CP2-size}
A category in Theorem~\ref{C:triangle-band-compression} with
$BC\simeq\mathbb CP^2$ satisfies $|M|+|S|\,|T|\geq23$.
In particular the fifteen-arrow profile
\[
 (|\End(0)|,|\End(1)|,|\End(2)|,
   |\Hom(0,1)|,|\Hom(1,2)|,|\Hom(0,2)|)
       =(1,8,1,2,2,1)
\]
is impossible: its compressed monoid would have only twelve
elements.
\end{corollary}
\begin{proof}
Apply Theorem~\ref{B:CP2-minimum23} to the compressed monoid.
\end{proof}
